Put \(D_{i,j}=\widehat\Psi_{i,j}-\Psi_j(O_i)\).  The score-expansion lemma gives
\(\max_j n^{-1}\sum_iD_{i,j}^2=o_P(1)\).  For fixed coordinates \(j,q\), the identity \(\widehat\Psi_{i,j}=\Psi_j(O_i)+D_{i,j}\) gives exactly
\[
 \widehat\Omega^{\mathrm{raw}}_{jq}-\Omega_{jq}
 =\left\{\frac1n\sum_i\Psi_j(O_i)\Psi_q(O_i)-\Omega_{jq}\right\}
  +\frac1n\sum_iD_{i,j}\widehat\Psi_{i,q}
  +\frac1n\sum_i\Psi_j(O_i)D_{i,q}.
\]
The fourth-moment assumptions and bounded propensity ratios imply \(E_P[|\Psi_j|^4]<\infty\).  Thus the first braced term is \(o_P(1)\), for example by applying Chebyshev's inequality to the i.i.d. products.  Moreover,
\[
 \frac1n\sum_i\widehat\Psi_{i,q}^2
 \le 2\frac1n\sum_i\Psi_q(O_i)^2+2\frac1n\sum_iD_{i,q}^2=O_P(1).
\]
Cauchy--Schwarz now bounds the last two terms by
\[
 \left(\frac1n\sum_iD_{i,j}^2\right)^{1/2}
 \left(\frac1n\sum_i\widehat\Psi_{i,q}^2\right)^{1/2}
 +
 \left(\frac1n\sum_i\Psi_j(O_i)^2\right)^{1/2}
 \left(\frac1n\sum_iD_{i,q}^2\right)^{1/2}=o_P(1).
\]
Because \(M\) is fixed, entrywise convergence implies
\(\|\widehat\Omega^{\mathrm{raw}}-\Omega\|_{\mathrm{op}}=o_P(1)\).

For clarity, nonexpansiveness of the positive-semidefinite projection needs no extra stochastic premise.  If \(C=\mathbb S_+^M\), the variational inequalities for Euclidean projection give
\(\langle A-\Pi_C A,Z-\Pi_C A\rangle_F\le0\) for every \(Z\in C\).  Applying this once with \(A=\widehat\Omega^{\mathrm{raw}}\), once with \(A=\Omega\), and adding proves
\[
 \|\Pi_C(\widehat\Omega^{\mathrm{raw}})-\Pi_C(\Omega)\|_F
 \le \|\widehat\Omega^{\mathrm{raw}}-\Omega\|_F.
\]
Since \(\Omega\in C\), \(\Pi_C(\Omega)=\Omega\), whence
\[
 \|\widehat\Omega-\Omega\|_{\mathrm{op}}=o_P(1).
\]
Finally,
\[
 \|\widehat\Omega_\lambda-\Omega\|_{\mathrm{op}}
 \le \|\widehat\Omega-\Omega\|_{\mathrm{op}}+\lambda_n=o_P(1),
\]
because \(\lambda_n=n^{-1/4}\to0\).